\begin{rubric}{Research Experience as PhD Candidate}

\parbox{\textwidth}{
  \textbf{Thesis title}: Complement Coercion and Logical Metonymy in Norwegian. A Computational Linguistic Research.\par
  % \textbf{Supervisors}: Giosuè Baggio, Alessandro Lenci, Emmanuele Chersoni.
  \textbf{Topics}: LLMs, Compositionality, Lexical Semantics, Commonsense Knowledge\par
% Research Methods: Multifaceted approach to investigate complement coercion and logical metonymy within Large Language Models (LLMs). My project studies their performance behavior towards phenomena containing compositionality violations, bringing to missing lexical and semantic information for a full understanding of the entire sentence meaning. Our aim is to understand how and to what extent LLMs can extrapolate underlying semantic and syntactic information for a prediction of plausible eventive interpretations, only by considering specific contextual information overtly expressed in textual form. The project will consider different linguistic features considered essential for solving complement coercion: syntactic structure, context in sentence, and lexical semantics of the arguments in coercion constituents.

\textbf{Research Methods}: Investigating how models handle compositionality violations and infer plausible eventive interpretations from limited contextual cues.

\textbf{Techincal Skills}: 
\begin{itemize}
    \item \textbf{Data}: Corpus analysis, data mining, data preprocessing and analysis, annotation.
    \item \textbf{Probing LLMs}: creation of datasets, Implementation of pretrained  autoregressive and autoencoder models (BERT, GPT-2, Llama, etc.), statistical analysis of performance and surprisal estimates.

\end{itemize}


% \textbf{Techincal Skills}

% \begin{itemize}
%     \item Data Collection and Data Mining
%     \begin{itemize}
%         \item Corpus analysis of large text corpora using NLP tools in Python (NLTK, Stanza, spaCy) for extraction of coercion sentences. Applied tokenization, lemmatization, stemming, POS tagging, and dependency parsing.
        
%         \item Preprocessing and storage of relevant sentences into structured files (e.g., CSV).
        
%         \item Classification and annotation of relevant sentences for quantitative and statistical analysis.
        
%         \item Experience in survey data collection using tools like Nettskjema.
%     \end{itemize}
%     \item Implementation of probing task for testing complement coercion in Norwegian

%         \begin{itemize}
%             \item Dataset creation with sentence template tailored for the task.
%             \item Familiarity with transformer-based LLMs in Norwegian with focus on both autoencoders and autoregressive frameworks. 
%             \item Task proposed: masked language modeling (cloze task) and next word prediction task.
%             \item Prediction annotation and performance evaluation. Metrics employed: mean average precision(mAP) and precision@1 (P@1).  
%             \item Further computation of surprisal estimates for the analysis of LLMs behavior towards different coercive and non-coercive syntactic structures. Regression models (e.g., statsmodel) were used as test statistics for detecting specific effect of syntactic constructions, lexical and verb types responsible for the influence of higher surprisal values. 
%         \end{itemize}
% \end{itemize}

% \textbf{Research Visiting}:
%     \begin{itemize}
%         \item Institute for Natural Language Processing (IMS). Head of the department: Sebastian Padó. April-June 2025. Exploration of Transformer variants (e.g., Siamese sentence-BERT models) and their performance evaluation on complement coercion and logical metonymy phenomena.
%         \item Computational Linguistics Laboratory at the University of Pisa (CoLing Lab), Italy. January-March 2024 (3 months). Head of the department: Alessandro Lenci. Exploration of new techniques for testing LLMs, especially considering lexical semantic and distributional semantics. 
%     \end{itemize}
    
\textbf{Extra Research networks I collaborated with}:
    \begin{itemize}
        \item Collaborative research among scholars from various institutions worldwide. Research head: Karen fort, Sorbonne Université / LORIA. Paper Work: Your Stereotypical Mileage May Vary: Practical Challenges of Evaluating Biases in Multiple Languages and Cultural Contexts. Link: https://inria.hal.science/hal-04537096/. Multilingual bias evaluation in masked language models (MLMs) using across different languages and cultural contexts.
    \end{itemize}

Timeline: 
End of contract in October 2025, expected date of PhD defense in February-March 2026. \textbf{Open to start a new position earlier}. 
}
\end{rubric}